\documentclass[../../../main.tex]{subfiles}
\begin{document}
\subsection{Identical Particle}
Indistinguishability implies that the Hamiltonian cannot change if the degrees of freedom of two arbitrary particles $i$ and $j$ are exchanged
\begin{equation*}
  \mathbb{H } = \mathbb{H} (\dots,i\dots,j,\dots) = \mathbb{H} (\dots,j\dots,i,\dots)
\end{equation*}
This can be seen by acting with the operator exchanging two arbitrary particles, $i$ and $j$, on a state of $N$ identical particles
\begin{equation*}
  \mathcal{T }_{ij} \ket{\dots, \phi_i,\dots,\phi_j,\dots} = \ket{\dots, \phi_j,\dots,\phi_i,\dots} 
\end{equation*}
which produces a state that may not be physically different from the original state, so that there can only be a phase factor $\nu_{ij }=e ^{i\alpha_{ij}}$
\begin{equation*}
  \mathcal{T }_{ij} \ket{\dots, \phi_i,\dots,\phi_j,\dots} = \nu_{ij }\ket{\dots, \phi_i,\dots,\phi_j,\dots} 
\end{equation*}
Moreover, because of $\mathcal{T}_{ij }=\mathcal{T}_{ji}$ we have
\begin{equation*}
  \mathcal{T}_{ij }^{2 }=1\qquad\nu_{ij }=\pm1
\end{equation*}
These results imply that there are only two types of quantum particles
\begin{enumerate}
  \item $\nu=+1$ particles whose states are symmetrical under particle exchange,
  \item $\nu=-1$ particles whose states are antisymmetrical under particle exchange.
\end{enumerate}

\subsection{Vector Potential}
The Coulomb gauge definition of vector potential $\nabla \cdot  \mathbf{A}=0$ constraint the vector to only the transversal (perpendicular) to the wave vector $\mathbf{k}$.
This can be seen by considering a complex monochromatic plane wave $\mathbf{A}(\mathbf{r},t)=\mathbf{A}_0 e^{i(\mathbf{k}\cdot \mathbf{r }-\omega t)}$.
Under Coulomb gauge
\begin{equation*}
  \nabla  \cdot \mathbf{A}=  \nabla \cdot  \mathbf{A}_0  e^{i(k \cdot \mathbf{r} - \omega t)} = \mathbf{A}_0 \cdot (i \mathbf{k}) e^{i(k \cdot \mathbf{r} - \omega t)} = \mathbf{k } \cdot \mathbf{A} = 0
\end{equation*}
This complex monochromatic plane wave is valid solution to wave equation that the vector potential obey.

We can expand the vector potential $\mathbf{A}(r, t)$ as a Fourier series of monochromatic plane waves.
\begin{equation*}
  \mathbf{A}(\mathbf{r}, t)
  = \sum_{\mathbf{k},s} A_{\mathbf{k} s}(t) \frac{e^{i \mathbf{k} \cdot \mathbf{r}}}{\sqrt{V}} \boldsymbol{\epsilon}_{\mathbf{k} s}
\end{equation*}
It follows from  classical electrodynamics that the vector potential is a real vector field $\mathbf{A} = \mathbf{A}^*$.
Choosing real polarization vectors $\boldsymbol{\epsilon}_{\mathbf{k} s}=\epsilon ^*_{\mathbf{k} s}$
\begin{equation*}
  \mathbf{A}(\mathbf{r},t) = \sum_{\mathbf{k},s} A^*_{\mathbf{k} s}(t)  \frac{e^{-i \mathbf{k} \cdot \mathbf{r}}}{\sqrt{V}}\boldsymbol{\epsilon}_{\mathbf{k} s}
  = \sum_{\mathbf{k},s} A^*_{-\mathbf{k} s}(t)  \frac{e^{i \mathbf{k} \cdot \mathbf{r}}}{\sqrt{V}}\boldsymbol{\epsilon}_{\mathbf{k} s}
\end{equation*}
where we relabel $\mathbf{k} \rightarrow -\mathbf{k}$.
Consequently $A_{-\mathbf{k},s}=A^*_{\mathbf{k},s}$.
Thus, we can sum the series over independent $\mathbf{k}$
\begin{equation*}
  \mathbf{A}(\mathbf{r}, t) = \sum_{\mathbf{k},s} \left[ A_{\mathbf{k}s}(t) \frac{e^{i\mathbf{k} \cdot \mathbf{r}}}{\sqrt{V}} + A_{\mathbf{k}s}^*(t) \frac{e^{-i\mathbf{k} \cdot \mathbf{r}}}{\sqrt{V}} \right] \boldsymbol{\epsilon}_{\mathbf{k}s}
\end{equation*}
where the complex plane waves $e^{i \mathbf{k } \cdot \mathbf{r}}$ normalized in a volume $V$ are the basis functions of the expansion, and $\boldsymbol{\epsilon}_{\mathbf{k}s}$ are mutually orthogonal real unit vectors of polarization which are also orthogonal to $\mathbf{k}$.
Note that the first representation count over all possible $\mathbf{k}$: both $\mathbf{k}$ and $-\mathbf{k}$ are included; meanwhile the second only the representation $\left\{ \mathbf{k},-\mathbf{k} \right\} $ are included.

Aside, it is called the Coulomb gauge because it would produce the Coulomb potential.
Start form the gauge $\nabla \cdot \mathbf{A}=0$, then consider the electric field enforced by the gauge
\begin{equation*}
  \mathbf{E}= -\nabla \phi+\frac{\partial \mathbf{A }}{\partial T}=-\nabla \phi
\end{equation*}
Then take the divergence
\begin{equation*}
  \nabla \cdot \mathbf{E }= -\nabla ^2 \phi
\end{equation*}
According to Gauss law
\begin{equation*}
  \nabla ^2 \phi=-\frac{\rho(\mathbf{r},t )}{\epsilon_0}
\end{equation*}
which is the Poisson equation whose solution is the Coulomb potential
\begin{equation*}
  \phi(\mathbf{r},t)=\frac{1}{4\pi \epsilon_0}\int\frac{1}{\rcurs}\rho(\mathbf{r}',t)d\tau'
\end{equation*}
where $\mathbf{r}$ is the field point (point in space where you are evaluating), $\mathbf{r}'$ is the source point (the location in space of an element of charge), and $\rcurs$ as the distance between source and field point.

\subsubsection{Wave equation.}
Each mode amplitude $A_{\mathbf{k}s}(t)$ behaves exactly like a classical harmonic oscillator due to $\mathbf{A}(\mathbf{r},t)$ obeying the wave equation that yield harmonic oscillator differential equation, whose solution is $A_{\mathbf{k}s}(t)=A_{\mathbf{k}s}(0)e^{i\omega_kt}$.
This can be proven by
\begin{align*}
  \left( \nabla^2 - \frac{1}{c^2} \frac{\partial^2}{\partial t^2} \right)    \sum_{\mathbf{k},s} A_{\mathbf{k} s}(t) \frac{e^{i \mathbf{k} \cdot \mathbf{r}}}{\sqrt{V}} \boldsymbol{\epsilon}_{\mathbf{k} s}
   & = 0  \\
  \sum_{\mathbf{k },s} \left[ -k^2 A_{\mathbf{k }s}(t) - \frac{1}{c^2} \frac{d^2 }{dt^2}A_{\mathbf{\mathbf{k}s}}(t) \right] \frac{e^{i \mathbf{k }\cdot \mathbf{r}}}{\sqrt{V} \epsilon_{\mathbf{k }s}}
   & =  0
\end{align*}
Each mode is independent, so
\begin{equation*}
  \frac{d^2 }{dt^2 }A_{\mathbf{k }s}(t)+\omega_k^2A_{\mathbf{A }s}(t)=0
\end{equation*}
with $\omega_k=ck$, which have the complex solutions
\begin{equation*}
  A_{\mathbf{k }s }(0)e^{-\omega_kt}
\end{equation*}

\subsection{Fock State}
The operators $\hat{a}_{\mathbf{k }s }$ and $\hat{a }_{\mathbf{k }s}$ act in Fock space $\mathcal{F}$.
A generic state of this Fock space F is given by
\begin{equation*}
  \ket{\left\{ n_{\mathbf{k }s} \right\} }
  = \ket{n_{\mathbf{k} s}n_{\mathbf{k}'s'}n_{\mathbf{k}''s''} \cdots }
  =\bigotimes_{\mathbf{k },s} \ket{n_{\mathbf{k }s}}
\end{equation*}
where there are \( n_{\mathbf{k} s} \) photons with wavevector \( \mathbf{k} \) and polarization \( s \), \( n_{\mathbf{k}'s'} \) photons with wavevector \( \mathbf{k}' \) and polarization \( s' \), \( n_{\mathbf{k}''s''} \) photons with wavevector \( \mathbf{k}'' \) and polarization \( s'' \), et cetera.
For example, the state $\ket{ 0_{\mathbf{k} s} 1_{\mathbf{k}'s'}  2_{\mathbf{k}''s''}}$.
Here, the entire mode $\mathbf{k},s$ is in vacuum or no photon exist; mode $\mathbf{k'},s'$ has one photon; and mode $\mathbf{k}'',s''$ has two.

The Fock space \( \mathcal{F} \) is given by
\begin{equation*}
  \mathcal{F} = \mathcal{H}_0 \oplus \mathcal{H}_1 \oplus \mathcal{H}_2 \oplus \mathcal{H}_3 \oplus \cdots = \bigoplus_{n=0}^{\infty} \mathcal{H}_n = \bigoplus_{n=0}^{\infty} \mathcal{H}^{\otimes n}
\end{equation*}
where $  \mathcal{H}_n = \mathcal{H} \otimes \mathcal{H} \otimes \cdots \otimes \mathcal{H} = \mathcal{H}^{\otimes n}$ is the Hilbert space of $n$ identical photons, which is $n$ times the tensor product of the single-photon Hilbert space $\mathcal{H}$
Notice that in the definition of the Fock space $\mathcal{F}$ one must include the space $\mathcal{H}_0 = \mathcal{H}^{\otimes 0}$, which is the Hilbert space of 0 photons, containing only the vacuum state
\begin{equation*}
  \ket{0}
  = \ket{0_{\mathbf{k} s}0_{\mathbf{k}'s'}\cdots }
  =\bigotimes_{\mathbf{k },s} \ket{0_{\mathbf{k }s}}
\end{equation*}

\subsubsection{Ladder operator.}
If there are exactly $n$ photons in this polarized monochromatic wave $\omega=\omega_k=c|\mathbf{k}|$ the Fock state of the system is given by
\begin{equation*}
  \ket{n }= \frac{1 }{\sqrt{n !}}\left( \hat{a }^\dagger \right) ^n \ket{0}
\end{equation*}
or more precisely
\begin{equation*}
  \ket{n_{\mathbf{k }s} }= \frac{1 }{\sqrt{n_{\mathbf{k }s}!}}\left( \hat{a }^\dagger_{\mathbf{k }s} \right) ^{n_{\mathbf{k }s}} \ket{0}
\end{equation*}
This can be generalized into  multimode state
\begin{equation*}
  \ket{\left\{ n_{\mathbf{k }s}  \right\} }
  =\prod_{\mathbf{k },s} \frac{1 }{\sqrt{n_{\mathbf{k }s}!}}\left( \hat{a }^\dagger_{\mathbf{k }s} \right) ^{n_{\mathbf{k }s}} \ket{0}
\end{equation*}

To prove this, we define unnormalized state
\begin{equation*}
  \ket{n'_{\mathbf{k }s}}=\hat{a}^\dagger_{\mathbf{k }s} \ket{0}
\end{equation*}
and use the relations
\begin{gather*}
  \hat{a}_{\mathbf{k }s} (\hat{a}^\dagger_{\mathbf{k }s} )^n
  = (\hat{a}^\dagger_{\mathbf{k }s} )^n a +n (\hat{a}^\dagger_{\mathbf{k }s} )^{n-1} \\
  \hat{a}_{\mathbf{k }s} ^n(\hat{a}^\dagger_{\mathbf{k }s}) ^n \ket{0}= n! \ket{0}
\end{gather*}

First let us use induction to prove said relation.
With both being the induction hypothesis, our base case $n$ is the commutator relation
\begin{equation*}
  \hat{a}_{\mathbf{k }s} \hat{a}^\dagger_{\mathbf{k }s} =\hat{a}^\dagger_{\mathbf{k }s} \hat{a}_{\mathbf{k }s} +1
\end{equation*}
We prove for $n+1$ by considering
\begin{align*}
  \hat{a}_{\mathbf{k }s} (\hat{a}^\dagger_{\mathbf{k }s} )^{n_{\mathbf{k }s}+1}
                                                                 &
  = \hat{a}_{\mathbf{k }s} (\hat{a}^\dagger_{\mathbf{k }s} )^{n_{\mathbf{k }s}} \hat{a}^\dagger_{\mathbf{k }s}
  =[(\hat{a}^\dagger_{\mathbf{k }s} )^{n_{\mathbf{k }s}} a +{n_{\mathbf{k }s}} (\hat{a}^\dagger_{\mathbf{k }s} )^{{n_{\mathbf{k }s}}-1}]\hat{a}^\dagger_{\mathbf{k }s}                                                    \\
                                                                 &
  = (\hat{a}^\dagger_{\mathbf{k }s} )^{n_{\mathbf{k }s}} \hat{a}_{\mathbf{k }s} \hat{a}^\dagger_{\mathbf{k }s}+ {n_{\mathbf{k }s}} (\hat{a}^\dagger_{\mathbf{k }s} )^{{n_{\mathbf{k }s}}-1}\hat{a}^\dagger_{\mathbf{k }s} \\
                                                                 &
  =  (\hat{a}^\dagger_{\mathbf{k }s} )^{n_{\mathbf{k }s}} (\hat{a}^\dagger_{\mathbf{k }s}  \hat{a}_{\mathbf{k }s} +1)+{n_{\mathbf{k }s}} (\hat{a}^\dagger_{\mathbf{k }s} )^{n_{\mathbf{k }s}}                             \\
  \hat{a}_{\mathbf{k }s} (\hat{a}^\dagger_{\mathbf{k }s} )^{n+1} & =
  (\hat{a}^\dagger_{\mathbf{k }s} )^{{n_{\mathbf{k }s}}+1}\hat{a}_{\mathbf{k }s} + ({n_{\mathbf{k }s}}+1)(\hat{a}^\dagger_{\mathbf{k }s} )^{n_{\mathbf{k }s}}
\end{align*}
which is our hypothesis for ${n_{\mathbf{k }s}} \to {n_{\mathbf{k }s}}+1$.
The action of this relation is as follows
\begin{equation*}
  \hat{a}_{\mathbf{k }s} (\hat{a}^\dagger_{\mathbf{k }s} )^{n_{\mathbf{k }s}} \ket{0}
  =(\hat{a}^\dagger_{\mathbf{k }s} )^{n_{\mathbf{k }s}} a +{n_{\mathbf{k }s}} (\hat{a}^\dagger_{\mathbf{k }s} )^{{n_{\mathbf{k }s}}-1}   \ket{0}= {n_{\mathbf{k }s}} (\hat{a}^\dagger_{\mathbf{k }s} )^{{n_{\mathbf{k }s}}-1}\ket{0}
\end{equation*}
Next we use this relation, which has been proven, to prove the second relation.
Using induction again, we consider base case $n$ as
\begin{equation*}
  (\hat{a}_{\mathbf{k }s} )^0 (\hat{a}^\dagger_{\mathbf{k }s} )^0 \ket{0} = \ket{0}=0!\ket{0}
\end{equation*}
It holds, then for $n+1$ consider
\begin{equation*}
  \hat{a}_{\mathbf{k }s} (\hat{a}^\dagger_{\mathbf{k }s} )^{{n_{\mathbf{k }s}}+1} \ket{0}
  = ({n_{\mathbf{k }s}}+1) (\hat{a}^\dagger_{\mathbf{k }s} )^{{n_{\mathbf{k }s}}}\ket{0}
\end{equation*}
then apply $\hat{a}_{\mathbf{k }s} ^{n_{\mathbf{k }s}}$ along with the second hypothesis
\begin{align*}
  (\hat{a}_{\mathbf{k }s})^{{n_{\mathbf{k }s}}+1} (\hat{a}^\dagger_{\mathbf{k }s} )^{{n_{\mathbf{k }s}}+1} \ket{0}
   & =  ({n_{\mathbf{k }s}}+1) (\hat{a}_{\mathbf{k }s} )^{n_{\mathbf{k }s}}(\hat{a}^\dagger_{\mathbf{k }s} )^{{n_{\mathbf{k }s}}}\ket{0} \\
   & = ({n_{\mathbf{k }s}}+1){n_{\mathbf{k }s}}! \ket{0}=({n_{\mathbf{k }s}}+1)\ket{0}
\end{align*}
It holds once more.
We now back at normalizing the state $\ket{{n_{\mathbf{k }s}}'}$ by taking the inner product and inserting the normalizing constant
\begin{equation*}
  \alpha^2\braket{{n_{\mathbf{k }s}}' | {n_{\mathbf{k }s}}'}=\alpha \braket{0 | \hat{a}_{\mathbf{k }s} ^{n_{\mathbf{k }s}} (\hat{a}^\dagger_{\mathbf{k }s} )^{n_{\mathbf{k }s}}|0}={n_{\mathbf{k }s}}!\\
  \alpha=\frac{1 }{\sqrt{{n_{\mathbf{k }s}}!}}
\end{equation*}
Thus the normalized state
\begin{equation*}
  \ket{{n_{\mathbf{k }s}}}=\alpha \ket{{n_{\mathbf{k }s}}'}= \frac{1 }{\sqrt{{n_{\mathbf{k }s}}! }}(\hat{a}^\dagger_{\mathbf{k }s}   )^{n_{\mathbf{k }s}} \ket{0}
\end{equation*}

For the multimode state, we can start by
\begin{equation*}
  \ket{\left\{ n_{\mathbf{k }s} \right\} } = \bigotimes_{\mathbf{k }s}\ket{n_{\mathbf{k }s}}=\bigotimes_{\mathbf{k }s}\frac{1 }{\sqrt{{n_{\mathbf{k }s}}! }}(\hat{a}^\dagger_{\mathbf{k }s}   )^{n_{\mathbf{k }s}} \ket{0}
\end{equation*}
Using
\begin{equation*}
  \bigotimes_i c_i \ket{n_i}= \left( \prod_i c_i  \right) \bigotimes_i \ket{n_i}
\end{equation*}
yields
\begin{equation*}
  \ket{\left\{ n_{\mathbf{k }s}  \right\} }
  =\prod_{\mathbf{k },s} \frac{1 }{\sqrt{n_{\mathbf{k }s}!}}\left( \hat{a }^\dagger_{\mathbf{k }s} \right) ^{n_{\mathbf{k }s}} \bigotimes_{\mathbf{k }s}\ket{0}
  =\prod_{\mathbf{k },s} \frac{1 }{\sqrt{n_{\mathbf{k }s}!}}\left( \hat{a }^\dagger_{\mathbf{k }s} \right) ^{n_{\mathbf{k }s}} \ket{0}
\end{equation*}

\subsubsection{(In)Coherent state.}
By using the number operator, one obtain
\begin{equation*}
  \hat{N}_{\mathbf{k }s} | \cdots n_{\mathbf{k }s} \cdots \rangle = n_{\mathbf{k }s}  | \cdots n_{\mathbf{k }s}  \cdots \rangle
\end{equation*}
Viewed in this light, the Fock state is the eigenstate of number operator and has a definite photon number, read: eigenvalue.

Despite this, the Fock state has no coherent phase and thus unable to properly interference.
This result can be seen by considering the expectation value of electromagnetic field.
Notice that
\begin{equation*}
  \hat{\mathbf{E}} \propto \sum_{\mathbf{k},s} \left( \hat{a}_{\mathbf{k }s} -\hat{a}^\dagger_{\mathbf{k }s}  \right) \quad
  \hat{\mathbf{B}} \propto \sum_{\mathbf{k},s} \left( \hat{a}_{\mathbf{k }s} -\hat{a}^\dagger_{\mathbf{k }s}  \right)
\end{equation*}
However, the operators change the occupation number of that mode and produce a state orthogonal to the original one
\begin{align*}
  \braket{\left\{ n_{\mathbf{k }s} \right\}   | \hat{a}_{\mathbf{k }s} |\left\{ n_{\mathbf{k }s} \right\} }
   & =   \sqrt{n_{\mathbf{k }s}} \braket{\left\{ n_{\mathbf{k }s} \right\}  | \left\{ n_{\mathbf{k }s} \right\} -1_{\mathbf{k }s}} =0     \\
  \braket{\left\{ n_{\mathbf{k }s} \right\}  | \hat{a}^\dagger_{\mathbf{k }s}  |\left\{ n_{\mathbf{k }s} \right\} }
   & =   \sqrt{n_{\mathbf{k }s}+1} \braket{\left\{ n_{\mathbf{k }s} \right\}  | \left\{ n_{\mathbf{k }s} \right\} + 1_{\mathbf{k }s}} = 0
\end{align*}
Thus
\begin{equation*}
  \braket{\left\{ n_{\mathbf{k }s} \right\}  | \hat{\mathbf{E }}|\left\{ n_{\mathbf{k }s} \right\} }=0\qquad
  \braket{\left\{ n_{\mathbf{k }s} \right\}  | \hat{\mathbf{B }}|\left\{ n_{\mathbf{k }s} \right\} }=0
\end{equation*}
The vanishing result does not mean the field is absent, it means that the average electromagnetic field amplitude is zero and field fluctuates symmetrically about zero.
The vanishing result thus can be seen as a consequence of undefined phase of Fock space.

The square of the field is different case, however.
With the fields written as
\begin{align*}
  \hat{\mathbf{E}}^2 & \propto \sum_{\mathbf{k},s}\sum_{\mathbf{k}',s'} \left( -\hat{a}_{\mathbf{k }s}\hat{a}_{\mathbf{k }'s'} +\hat{a}_{\mathbf{k }s} \hat{a}^\dagger_{\mathbf{k }'s'} +\hat{a}^\dagger_{\mathbf{k }s} \hat{a}_{\mathbf{k }'s'}  -\hat{a}^\dagger_{\mathbf{k }s} \hat{a}^\dagger_{\mathbf{k }'s'}  \right) \\
  \hat{\mathbf{B}}^2 & \propto \sum_{\mathbf{k},s}\sum_{\mathbf{k}',s'} \left( -\hat{a}_{\mathbf{k }s}\hat{a}_{\mathbf{k }'s'} +\hat{a}_{\mathbf{k }s} \hat{a}^\dagger_{\mathbf{k }'s'} +\hat{a}^\dagger_{\mathbf{k }s} \hat{a}_{\mathbf{k }'s'}  -\hat{a}^\dagger_{\mathbf{k }s} \hat{a}^\dagger_{\mathbf{k }'s'}  \right)
\end{align*}
where we have inserted the factor $i$.
Along with
\begin{gather*}
  \sum_{\mathbf{k},s}\sum_{\mathbf{k}',s'} \braket{\hat{a}_{\mathbf{k }s}   \hat{a}_{\mathbf{k }'s'} }=0                        \\
  \sum_{\mathbf{k},s}\sum_{\mathbf{k}',s'} \braket{\hat{a}_{\mathbf{k }s}   \hat{a}^\dagger_{\mathbf{k }'s'} }=\sum_{\mathbf{k},s}\sum_{\mathbf{k}',s'} \delta_{\mathbf{k}\mathbf{k'}}\delta_{ss'}(n_{\mathbf{k }s}+1)= \sum_{\mathbf{k},s}(n_{\mathbf{k }s}+1) \\
  \sum_{\mathbf{k},s}\sum_{\mathbf{k}',s'} \braket{\hat{a}^\dagger_{\mathbf{k }s}=   \hat{a}_{\mathbf{k }'s'} } \sum_{\mathbf{k},s}\sum_{\mathbf{k}',s'} \delta_{\mathbf{k}\mathbf{k'}}\delta_{ss'}n_{\mathbf{k }s}=\sum_{\mathbf{k},s} n_{\mathbf{k }s} \\
  \sum_{\mathbf{k},s}\sum_{\mathbf{k}',s'} \braket{\hat{a}^\dagger_{\mathbf{k }s}   \hat{a}^\dagger_{\mathbf{k }'s'} }=0
\end{gather*}
the expectation value now can be evaluated as
\begin{equation*}
  \braket{\left\{ n_{\mathbf{k }s} \right\}  | \hat{\mathbf{E }}^2|\left\{ n_{\mathbf{k }s} \right\} }\propto \sum_{\mathbf{k},s}(2 n_{\mathbf{k }s}+1)
  \qquad
  \braket{\left\{ n_{\mathbf{k }s} \right\}  | \hat{\mathbf{B }}^2|\left\{ n_{\mathbf{k }s} \right\} }\propto \sum_{\mathbf{k},s} (2 n_{\mathbf{k }s}+1)
\end{equation*}
or with all terms inserted
\begin{align*}
  \braket{\left\{ n_{\mathbf{k }s} \right\}  | \hat{\mathbf{E }}^2|\left\{ n_{\mathbf{k }s} \right\} } & =  \sum_{\mathbf{k},s} \frac{\hbar \omega_k }{\epsilon_0 V}\left( 2n_{\mathbf{k }s}+1\right)
  \\
  \braket{\left\{ n_{\mathbf{k }s} \right\}  | \hat{\mathbf{B }}^2|\left\{ n_{\mathbf{k }s} \right\} } & =  \sum_{\mathbf{k},s} \frac{\hbar \omega_k \mu_0}{V}\left( 2n_{\mathbf{k }s}+1 \right)
\end{align*}
\begin{equation*}
\end{equation*}
We can insert this into classical Hamiltonian
\begin{align*}
  H & =   \int d^3 \mathbf{r}  \left( \frac{\epsilon_0}{2} \mathbf{E}^2 + \frac{1}{2\mu_0} \mathbf{B}^2 \right)        \\
    & = \int d^3 \mathbf{r } \sum_{\mathbf{k},s} \left[ \frac{\hbar  \omega_k }{V }\left( n_{\mathbf{k }s}+\frac{1 }{2 } \right)  \right] \\
  H & = \sum_{\mathbf{k},s} \hbar \omega_k \left( n_{\mathbf{k }s}+\frac{1 }{2} \right)
\end{align*}
This result match the quantum Hamiltonian
\begin{align*}
  \braket{\left\{ n_{\mathbf{k }s} \right\} |H|\left\{ n_{\mathbf{k }s} \right\} }
   & =
  \left\langle \left\{ n_{\mathbf{k }s} \right\} \Bigg| \sum_{\mathbf{k},s} \hbar \omega_k \left( \hat{N }_{\mathbf{k} s }+ \frac{1 }{2} \right)\Bigg|\left\{ n_{\mathbf{k }s} \right\}  \right\rangle \\
  \braket{H}
   & =  \sum_{\mathbf{k},s} \hbar \omega_k \left( n_{\mathbf{k }s} +\frac{1 }{2} \right)
\end{align*}
\begin{equation*}
\end{equation*}
This mean that the field have definite energy despite not having definite phase.
The fact that the energy is definite despite the field fluctuates is because both energy and phase are conjugate variables in the quantum.
This field fluctuation reflects uncertainty in phase.
\end{document}